% Readout 6: encoder + per-candidate / per-label scoring.
% (a) uni-encoder (GLiClass-style): labels with marker tokens + SEP + text in one
%     bidirectional sequence; label markers gathered, text mean-pooled, scored.
% (b) bi-encoder: text tower u:[B,d], label tower V:[K,d] (cacheable), z=uV^T/tau.
% Cross-encoder (one (text, candidate) row each) = readout 3 with an encoder.
% MetisCraft readout tensor figures: shared notation, palette and macros.
% Cell counts are illustrative; each symbolic axis keeps one physical edge
% length inside a figure (u per cell): T=12, d=5, V=12, K=L=3, B=2, d_h=4.
\documentclass[tikz,border=3mm]{standalone}
\usepackage[UTF8,fontset=fandol]{ctex}
\usepackage{amsmath,amssymb,bm}
\usetikzlibrary{arrows.meta,calc,backgrounds,decorations.pathreplacing}
\definecolor{Ink}{HTML}{273648}
\definecolor{Muted}{HTML}{607086}
% token roles
\definecolor{Ctx}{HTML}{85898F}    % state / context
\definecolor{Qst}{HTML}{4F8FA5}    % question
\definecolor{Opt}{HTML}{5A9A6E}    % options / candidates (three lightness levels)
\definecolor{Gen}{HTML}{B89A5E}    % generated rationale tokens
% tensor roles
\definecolor{Hid}{HTML}{8A74B5}    % hidden activations
\definecolor{Wgt}{HTML}{E0904F}    % trainable / pretrained weights
\definecolor{Lgt}{HTML}{C95B5B}    % logits, probabilities, loss
\definecolor{Idx}{HTML}{7E8B9B}    % integer indices / labels
\definecolor{Msk}{HTML}{3F5A78}    % binary attention support
\colorlet{TokS}{Ctx!55}\colorlet{TokQ}{Qst!65}
\colorlet{Tok1}{Opt!30}\colorlet{Tok2}{Opt!58}\colorlet{Tok3}{Opt!88}
\colorlet{TokD}{Ink!80}\colorlet{TokG}{Gen!65}\colorlet{TokM}{Opt!100!black!60}
\colorlet{TokC}{Ink!80}\colorlet{TokP}{white}
\colorlet{TokL}{Opt!58}\colorlet{TokT}{Qst!65}
\newcommand{\U}{0.3}
\newcommand{\Off}{0.14}
\newcounter{tseed}
\tikzset{
  flow/.style={-{Latex[length=2mm,width=1.4mm]},draw=Ink!85,line width=.7pt},
  dflow/.style={-{Latex[length=2mm,width=1.4mm]},draw=Ink!60,line width=.6pt,dashed},
  op/.style={draw=Ink!30,fill=white,rounded corners=2pt,line width=.5pt,
             align=center,font=\small,inner sep=4pt},
  sym/.style={font=\small,inner sep=1pt},
  shp/.style={font=\scriptsize,text=Muted,inner sep=1pt},
  note/.style={font=\scriptsize,text=Muted,align=center,inner sep=1pt},
  stage/.style={font=\small\bfseries,text=Muted},
  bigop/.style={font=\Large,text=Ink},
  brc/.style={decorate,decoration={brace,amplitude=3pt},draw=Ink!55,line width=.5pt},
}
% \Tens{id}{cx}{cy}{rows}{cols}{color}{depth}{symbol}{shape}
\newcommand{\Tens}[9]{%
  \stepcounter{tseed}%
  \pgfmathsetmacro{\tw}{#5*\U}\pgfmathsetmacro{\th}{#4*\U}%
  \pgfmathsetmacro{\tx}{#2-\tw/2}\pgfmathsetmacro{\ty}{#3-\th/2}%
  \ifnum#7>0
    \foreach \k in {#7,...,1}{%
      \path[draw=#6!45,fill=white,rounded corners=.8pt,line width=.4pt]
        ({\tx+\k*\Off},{\ty+\k*\Off}) rectangle ++(\tw,\th);}%
  \fi
  \foreach \r in {1,...,#4}{\foreach \c in {1,...,#5}{%
    \pgfmathtruncatemacro{\lv}{mod(int(abs(sin(\r*127.1+\c*311.7+\arabic{tseed}*74.7))*1000),7)}%
    \pgfmathtruncatemacro{\tone}{\lv<3?30:(\lv<6?55:80)}%
    \path[fill=#6!\tone,rounded corners=.6pt]
      ({\tx+(\c-1)*\U+.022},{\ty+\th-\r*\U+.022}) rectangle ++({\U-.044},{\U-.044});}}%
  \coordinate (#1-w) at (\tx,#3);
  \coordinate (#1-e) at ({\tx+\tw+#7*\Off},#3);
  \coordinate (#1-n) at (#2,{\ty+\th+#7*\Off});
  \coordinate (#1-s) at (#2,\ty);
  \coordinate (#1-nw) at (\tx,{\ty+\th});
  \coordinate (#1-sw) at (\tx,\ty);
  \node[sym] (#1-sym) at (#2,{\ty-.27}) {$#8$};
  \node[shp] (#1-shp) at (#2,{\ty-.6}) {$#9$};}
% outline one row (r, 1-based) of a tensor whose front face top-left is (x,y)
% \RowMark{x_left}{y_top}{row}{cols}{color}
\newcommand{\RowMark}[5]{%
  \path[draw=#5,line width=1.1pt,rounded corners=1pt]
    ({#1-.03},{#2-#3*\U-.03}) rectangle ++({#4*\U+.06},{\U+.06});}
\newcommand{\ColMark}[5]{% x_left y_top col rows color
  \path[draw=#5,line width=1.1pt,rounded corners=1pt]
    ({#1+(#3-1)*\U-.03},{#2-#4*\U-.03}) rectangle ++({\U+.06},{#4*\U+.06});}
% \TokRow{x_left}{y_center}{roles}: one tile per token, roles in {S,Q,1,2,3,D,G,M,C,P,L,T}
\newcommand{\TokRow}[3]{%
  \foreach \t [count=\i] in {#3}{%
    \path[fill=Tok\t,draw=Ink!25,line width=.3pt,rounded corners=.6pt]
      ({#1+(\i-1)*\U+.022},{#2-\U/2+.022}) rectangle ++({\U-.044},{\U-.044});}}
\newcommand{\TokCol}[3]{% x_center y_top roles (top to bottom)
  \foreach \t [count=\i] in {#3}{%
    \path[fill=Tok\t,draw=Ink!25,line width=.3pt,rounded corners=.6pt]
      ({#1-\U/2+.022},{#2-\i*\U+.022}) rectangle ++({\U-.044},{\U-.044});}}
% \Mask{x_left}{y_top}{n}{groups}{mode}
%  mode 0: causal & (g_j=0 or g_j=g_i)      shared-prefix tree
%  mode 1: all ones                          bidirectional
%  mode 2: causal & (g_j=0 or g_j=g_i or g_i=9)  pointer (9 = decision token)
%  mode 3: causal                            plain causal
%  mode 4: g_j=g_i                           independent bidirectional blocks
%  mode 5: causal with generated tokens (g=5) drawn in Gen colour
\newcommand{\Mask}[5]{%
  \foreach \i in {1,...,#3}{\foreach \j in {1,...,#3}{%
    \pgfmathtruncatemacro{\gi}{{#4}[\i-1]}%
    \pgfmathtruncatemacro{\gj}{{#4}[\j-1]}%
    \ifcase#5
      \pgfmathtruncatemacro{\ok}{and(\j<=\i,or(\gj==0,\gj==\gi))}\or
      \pgfmathtruncatemacro{\ok}{1}\or
      \pgfmathtruncatemacro{\ok}{and(\j<=\i,or(or(\gj==0,\gj==\gi),\gi==9))}\or
      \pgfmathtruncatemacro{\ok}{\j<=\i}\or
      \pgfmathtruncatemacro{\ok}{\gj==\gi}\or
      \pgfmathtruncatemacro{\ok}{\j<=\i}\fi
    \ifnum\ok=1
      \path[fill=Msk!70,rounded corners=.5pt]
        ({#1+(\j-1)*\U+.022},{#2-\i*\U+.022}) rectangle ++({\U-.044},{\U-.044});
    \fi}}%
  \path[draw=Ink!40,line width=.4pt] (#1,#2) rectangle ++({#3*\U},{-#3*\U});}
% integer slot strip: \Slots{x_left}{y_center}{entries}
\newcommand{\Slots}[3]{%
  \foreach \v [count=\i] in {#3}{%
    \node[draw=Idx!70,fill=Idx!12,rounded corners=1pt,minimum width=\U cm-.04cm,
      minimum height=\U cm-.04cm,inner sep=0pt,font=\tiny,text=Ink]
      at ({#1+(\i-.5)*\U},#2) {$\v$};}}
% bottom meaning box: \Meaning{x0}{y_top}{width}{axes}{objects}{mechanism}{id line}
\newcommand{\Meaning}[7]{%
  \node[anchor=north west,fill=Ink!4,rounded corners=3pt,inner sep=9pt,
    text width=#3cm-18pt,align=left,font=\small] (meaning) at (#1,#2) {%
    \begin{tabular}{@{}p{1.1cm}@{\hspace{6pt}}p{\dimexpr #3cm-2.0cm\relax}@{}}
    \textbf{轴} & \raggedright #4\tabularnewline[5pt]
    \textbf{对象} & \raggedright #5\tabularnewline[5pt]
    \textbf{机制} & \raggedright #6\tabularnewline
    \end{tabular}};
  \node[font=\scriptsize,text=Ink!45,anchor=north] at ($(meaning.south)+(0,-.18)$) {#7};}

\begin{document}
\begin{tikzpicture}[x=1cm,y=1cm,text=Ink]
\node[font=\large] at (11.0,4.05)
  {$e_{b,k}=H_{b,m_{b,k},:},\ \ u_b=\operatorname{mean}_{t\in\mathcal T_b}H_{b,t,:},\ \
    z_{b,k}=\langle e_{b,k},u_b\rangle\ \text{（同编码）};\qquad
    z=UV^{\mathsf T}/\tau\ \text{（双塔）};\qquad p_{b,k}=\sigma(z_{b,k})$};
% ---------- stage A: uni-encoder ----------
\node[stage] at (11.0,3.35) {(a) 同编码：标签与文本拼成一条序列，双向注意力，所有 token 互相可见};
\node[sym,anchor=east] at (1.0,2.3) {$X$};
\node[shp,anchor=east] at (1.0,2.0) {$[B,T]$};
\TokRow{1.2}{2.15}{M,1,M,2,M,3,C,T,T,T,T,T}
\draw[brc] (1.22,2.36) -- (2.98,2.36) node[midway,above=3pt,note] {标签 $\times K$};
\draw[brc] (3.32,2.36) -- (4.78,2.36) node[midway,above=3pt,note] {文本};
\node[note,anchor=north] at (3.15,1.98) {};
\Mask{1.2}{1.8}{12}{0,0,0,0,0,0,0,0,0,0,0,0}{1}
\node[sym] at (3.0,-2.07) {$M^{\mathrm{bi}}$};
\node[shp] at (3.0,-2.4) {$[T,T]$ 全 1（padding 除外）};
\node[op,minimum width=2.6cm,minimum height=1.3cm] (enc) at (6.75,0)
  {\textbf{Encoder}\\[2pt]双向注意力};
\draw[flow] (4.95,0) -- (enc.west);
\Tens{H}{9.35}{0}{12}{5}{Hid}{1}{H}{[B,T,d]}
\foreach \r in {1,3,5}{\RowMark{8.6}{1.8}{\r}{5}{TokM}}
\path[draw=Qst,line width=1.1pt,rounded corners=1pt] (8.57,-1.83) rectangle (10.13,-0.27);
\draw[flow] (enc.east) -- (H-w);
\node[op,minimum width=2.3cm,minimum height=1.0cm] (go) at (12.4,1.3)
  {\textbf{Gather 标记}\\[2pt]$H_{b,m_{b,k},:}$};
\node[op,minimum width=2.3cm,minimum height=1.0cm] (gp) at (12.4,-1.3)
  {\textbf{mean-pool}\\[2pt]文本位置 $\mathcal T_b$};
\draw[flow] (H-e) -- (go.west);
\draw[flow] (H-e) -- (gp.west);
\Slots{11.95}{2.25}{0,2,4}
\node[sym,anchor=west] at (12.9,2.25) {$m$\;{\scriptsize\color{Muted}$[B,K]$}};
\draw[dflow] (12.4,2.08) -- (go.north);
\Tens{E}{15.4}{1.3}{3}{5}{Hid}{1}{E}{[B,K,d]}
\Tens{u}{15.4}{-1.3}{2}{5}{Hid}{0}{u}{[B,d]}
\draw[flow] (go.east) -- (E-w);
\draw[flow] (gp.east) -- (u-w);
\node[op,minimum width=2.4cm,minimum height=1.3cm] (sc) at (18.6,0)
  {\textbf{scorer}\\[2pt]$\langle e_{b,k},u_b\rangle$\\[2pt]（或 MLP 等）};
\draw[flow] (E-e) -- (sc.north west);
\draw[flow] (u-e) -- (sc.south west);
\Tens{za}{21.25}{0}{2}{3}{Lgt}{0}{z}{[B,K]}
\draw[flow] (sc.east) -- (za-w);
% ---------- stage B: bi-encoder ----------
\node[stage] at (11.0,-3.45) {(b) 双塔：文本与标签各自编码，标签向量可离线缓存，候选之间互不可见};
\node[sym,anchor=east] at (1.0,-4.85) {$X^{\mathrm{txt}}$};
\node[shp,anchor=east] at (1.0,-5.15) {$[B,T_x]$};
\TokRow{1.2}{-5.0}{T,T,T,T,T,T,T,T}
\node[op,minimum width=2.4cm,minimum height=1.0cm] (et) at (5.25,-5.0)
  {Encoder + pool};
\draw[flow] (3.55,-5.0) -- (et.west);
\Tens{ub}{8.35}{-5.0}{2}{5}{Hid}{0}{u}{[B,d]}
\draw[flow] (et.east) -- (ub-w);
\node[sym,anchor=east] at (1.0,-7.25) {$X^{\mathrm{lab}}$};
\node[shp,anchor=east] at (1.0,-7.55) {$[K,T_\ell]$};
\TokRow{1.2}{-7.1}{C,1,1}
\TokRow{1.2}{-7.4}{C,2,P}
\TokRow{1.2}{-7.7}{C,3,3}
\node[op,minimum width=2.4cm,minimum height=1.0cm] (el) at (5.25,-7.4)
  {Encoder + pool\\[2pt]{\scriptsize 可离线缓存}};
\draw[flow] (2.25,-7.4) -- (el.west);
\Tens{V}{8.35}{-7.4}{3}{5}{Hid}{0}{V}{[K,d]}
\draw[flow] (el.east) -- (V-w);
\node[op,minimum width=2.2cm,minimum height=1.1cm] (mm) at (11.6,-6.2)
  {$uV^{\mathsf T}/\tau$\\[2pt]{\scriptsize $(B{\times}d)(d{\times}K)$}};
\draw[flow] (ub-e) -- (mm.north west);
\draw[flow] (V-e) -- (mm.south west);
\Tens{zb}{14.3}{-6.2}{2}{3}{Lgt}{0}{z}{[B,K]}
\draw[flow] (mm.east) -- (zb-w);
% ---------- stage C: normalisation and loss ----------
\node[stage] at (11.0,-9.3) {(a)(b) 共用：逐标签 sigmoid（多标签自然），单选也可沿 $K$ softmax};
\Tens{z}{1.95}{-10.8}{2}{3}{Lgt}{0}{z}{[B,K]}
\node[op,minimum width=2.4cm,minimum height=1.1cm] (sg) at (4.6,-10.8)
  {逐个 sigmoid\\[2pt]{\scriptsize 或 softmax 沿 $K$}};
\draw[flow] (z-e) -- (sg.west);
\Tens{p}{7.25}{-10.8}{2}{3}{Lgt}{0}{p}{[B,K]}
\draw[flow] (sg.east) -- (p-w);
\node[op,minimum width=2.4cm,minimum height=1.1cm] (loss) at (9.9,-10.8)
  {\textbf{BCE}（多标签）\\[2pt]或交叉熵};
\draw[flow] (p-e) -- (loss.west);
\Slots{9.45}{-11.9}{1,0,1}
\Slots{9.45}{-12.2}{0,0,1}
\node[sym,anchor=west] at (10.45,-12.05) {$Y$\;{\scriptsize\color{Muted}$\{0,1\}^{B\times K}$}};
\draw[dflow] (9.9,-11.73) -- (loss.south);
\node[note,anchor=west,align=left] at (12.4,-10.8)
  {交叉编码（Kenning v0.5、Luce \texttt{cross}）：每个 (文本, 候选) 拼成一行、行内双向，\\[3pt]
   读 [CLS] 或 pooling 得标量 —— 即读出 3把 decoder 换成 encoder，代价 $B\cdot K$ 行。};
% ---------- meaning ----------
\Meaning{0.0}{-12.85}{22.6}
 {$B$ 问题（文本）数；$K$ 标签 / 候选数；$T$ 同编码序列长度；$T_x,T_\ell$ 双塔中文本、标签长度；$d$ hidden；$\tau$ 温度。}
 {$m_{b,k}$ 第 $k$ 个标签标记 token 的下标；$\mathcal T_b$ 文本 token 的位置集合；$E$ 标签表示；$u$ 文本表示；$V$ 缓存的标签向量；$Y$ 多标签 0/1 目标。}
 {同编码时标签之间、标签与文本之间全可见（标签顺序会影响结果，训练常做标签打乱）；双塔时互不可见，标签可增删而不重算文本。encoder 小而快、CPU 可用，但长 state 与算术、规则类任务容量受限。}
 {读出 6：encoder + 逐候选 / 标签打分@Shichao}
\end{tikzpicture}
\end{document}
